\subsection{Smith Normal Form}

\begin{definition} \label{def:standing-pid-notation} \hlt{Standing PID Notation}\\
In the PID-structure results, $R$ denotes a principal ideal domain unless otherwise stated.\\
For $r\in R$, write $\langle r\rangle=Rr$; for an $R$-module $M$ and $v\in M$, write $\langle\!\langle v\rangle\!\rangle_R=Rv$.
\end{definition}

\begin{perspective} \label{per:pid-modules-scalar-arithmetic-controls-structure} \hlt{PID Arithmetic Controls Module Structure}\\
Fields have $F\setminus\{0\}\subseteq F^\times$. PIDs replace invertibility by factorization, gcds, and lcms.\\
This divisibility data classifies finitely generated modules.
\end{perspective}

\begin{definition} \label{def:module-order-over-pid} \hlt{Order of a Submodule or Element}\\
Let $M$ be an $R$-module and $N\leq M$. Since $\operatorname{Ann}_R(N)=\{r\in R:rN=0\}=\langle\alpha\rangle$, any such $\alpha$ is an \hlt{order} of $N$; write $o(N)=\alpha$, up to associates. For $v\in M$, set $o(v)\defeq o(\langle\!\langle v\rangle\!\rangle_R)$.
\end{definition}

\begin{proposition} \label{prop:orders-direct-sum-lcm} \hlt{Orders and Direct Sums}\\
Let $M=A\oplus B$. If $o(A)=\alpha$ and $o(B)=\beta$, then $o(M)\sim\operatorname{lcm}(\alpha,\beta)$.\\
More generally, $o(\bigoplus\limits_{i=1}^{n}M_i)\sim\operatorname{lcm}(o(M_1),\ldots,o(M_n))$.
\begin{proof}
$r(A\oplus B)=0\Longleftrightarrow rA=0$ and $rB=0$.\\
Hence $\operatorname{Ann}_R(A\oplus B)=\operatorname{Ann}_R(A)\cap\operatorname{Ann}_R(B) =\langle\alpha\rangle\cap\langle\beta\rangle =\langle\operatorname{lcm}(\alpha,\beta)\rangle$.\\
Thus $o(A\oplus B)\sim\operatorname{lcm}(\alpha,\beta)$, and induction gives the finite case.\\
\hlt{Structural Commentary}: Direct sums intersect annihilator conditions.\\
\hlt{Key Insight}: Lcms appear because one scalar must kill every summand at once.
\end{proof}
\end{proposition}

\begin{theorem} \label{thm:cyclic-modules-over-pid} \hlt{Cyclic Modules over a PID}\\
Let $M$ be an $R$-module and let $v\in M$.
\begin{enumerate}[label=(\roman*)]
\setlength{\itemsep}{0pt}
\item If $\operatorname{Ann}_R(v)=\langle\alpha\rangle$, then $R/\langle\alpha\rangle\xrightarrow{\sim}\langle\!\langle v\rangle\!\rangle_R$, $r+\langle\alpha\rangle\mapsto rv$.
\item If $C=\langle\!\langle v\rangle\!\rangle_R$ and $S\leq C$, then $S=\langle\!\langle sv\rangle\!\rangle_R$ for some $s\in R$.
\item If $o(v)=\alpha$ and $\beta\in R$, then $o(\beta v)\sim\alpha/\gcd(\alpha,\beta)$, and $\langle\!\langle\beta v\rangle\!\rangle_R=\langle\!\langle v\rangle\!\rangle_R \Longleftrightarrow \gcd(\alpha,\beta)\sim1$.
\end{enumerate}
\begin{proof}
\begin{enumerate}[label=(\roman*)]
\setlength{\itemsep}{0pt}
\item The map $\phi:R\to\langle\!\langle v\rangle\!\rangle_R$, $\phi(r)=rv$, is onto and $\ker\phi=\operatorname{Ann}_R(v)=\langle\alpha\rangle$. Hence $R/\langle\alpha\rangle\xrightarrow{\sim}\langle\!\langle v\rangle\!\rangle_R$.
\item For $S\leq C$, set $I=\{r\in R:rv\in S\}\lhd R$.\\
Since $R$ is a PID, $I=\langle s\rangle$ for some $s\in R$, and $S=\{rv:r\in I\}=\langle\!\langle sv\rangle\!\rangle_R$.
\item Put $d=\gcd(\alpha,\beta)$, $\alpha=d\alpha'$, and $\beta=d\beta'$ with $\gcd(\alpha',\beta')\sim1$.\\
For $r\in R$, $r\beta v=0\Longleftrightarrow\alpha\mid r\beta \Longleftrightarrow\alpha'\mid r\beta'\Longleftrightarrow\alpha'\mid r$, the last step by Bezout.\\
Thus $\operatorname{Ann}_R(\beta v)=\langle\alpha'\rangle =\langle\alpha/\gcd(\alpha,\beta)\rangle$. Finally,
$\langle\!\langle\beta v\rangle\!\rangle_R=\langle\!\langle v\rangle\!\rangle_R$ $\Longleftrightarrow(\exists c\in R)(v=c\beta v)$\\
$\Longleftrightarrow(\exists c\in R)(\alpha\mid1-c\beta) \Longleftrightarrow\gcd(\alpha,\beta)\sim1$.
\end{enumerate}
\noindent
\hlt{Structural Commentary}: A cyclic module is the base ring modulo the first relation satisfied by its generator.\\
\hlt{Key Insight}: Multiplying a cyclic generator divides its order by the gcd of the multiplier and the old order.
\end{proof}
\end{theorem}

\begin{theorem} \label{thm:cyclic-composition-coprime-orders} \hlt{Composing and Decomposing Cyclic Modules}\\
Let $M$ be an $R$-module.
\begin{enumerate}[label=(\roman*)]
\setlength{\itemsep}{0pt}
\item If $o(u_i)=\alpha_i\neq0$ and $\gcd(\alpha_i,\alpha_j)\sim1$ for $i\neq j$, then $o(\sum\limits_{i=1}^{n}u_i)\sim\prod_{i=1}^{n}\alpha_i$ and $\bigoplus\limits_{i=1}^{n}\langle\!\langle u_i\rangle\!\rangle_R =\langle\!\langle\sum\limits_{i=1}^{n}u_i\rangle\!\rangle_R$.
\item If $o(v)=\alpha=\prod\limits_{i=1}^{n}\alpha_i\neq0$ and $\gcd(\alpha_i,\alpha_j)\sim1$ for $i\neq j$,\\
then $\exists v_1,\ldots,v_n\in M$ such that $v=\sum\limits_{i=1}^{n}v_i$, $o(v_i)\sim\alpha_i$, and $\langle\!\langle v\rangle\!\rangle_R =\bigoplus\limits_{i=1}^{n}\langle\!\langle v_i\rangle\!\rangle_R$.
\end{enumerate}
\begin{proof}
\begin{enumerate}[label=(\roman*)]
\setlength{\itemsep}{0pt}
\item It suffices to prove $n=2$ and iterate. Let $o(u)=\alpha$, $o(w)=\beta$, and $\gcd(\alpha,\beta)\sim1$.\\
If $x\in\langle\!\langle u\rangle\!\rangle_R\cap\langle\!\langle w\rangle\!\rangle_R$, then $\alpha x=\beta x=0$, so Bezout gives $x=0$.  If $r(u+w)=0$, then $ru=-rw$ lies in this intersection; hence $ru=rw=0$, so $\alpha\mid r$ and $\beta\mid r$, hence $\alpha\beta\mid r$. Since $\alpha\beta(u+w)=0$, $o(u+w)\sim\alpha\beta$. Choose $a,b\in R$ with $a\alpha+b\beta=1$. Then $u=b\beta(u+w)$ and $w=a\alpha(u+w)$, so $\langle\!\langle u\rangle\!\rangle_R\oplus \langle\!\langle w\rangle\!\rangle_R=\langle\!\langle u+w\rangle\!\rangle_R$.
\item Put $\beta_i=\alpha/\alpha_i$. Since $\gcd(\beta_1,\ldots,\beta_n)\sim1$, choose $c_i\in R$ with $\sum\limits_{i=1}^{n}c_i\beta_i=1$, and set $v_i=c_i\beta_i v$.\\
Then $v=\sum\limits_{i=1}^{n}v_i$. Also $\gcd(c_i,\alpha_i)\sim1$ and $\gcd(\alpha,\beta_i)\sim\beta_i$, so Theorem \ref{thm:cyclic-modules-over-pid} gives $o(v_i)\sim\alpha_i$.\\
Part (i) gives $\langle\!\langle v\rangle\!\rangle_R =\bigoplus\limits_{i=1}^{n}\langle\!\langle v_i\rangle\!\rangle_R$.
\end{enumerate}
\noindent
\hlt{Structural Commentary}: Relatively prime cyclic pieces merge into one cyclic piece.\\
\hlt{Key Insight}: Bezout identities recover each component from the single sum.
\end{proof}
\end{theorem}

\begin{perspective} \label{per:pid-noetherian-specialization} \hlt{Noetherian Modules over a PID}\\
Over a PID, $N\leq R^n$ implies $N$ free with $\operatorname{rk}_R(N)\leq n$; Noetherianity is the f.g. shadow.
\end{perspective}

\begin{theorem} \label{thm:submodule-of-finite-free-pid-is-free} \hlt{Submodules of Finite Free PID Modules}\\
Let $F$ be free over $R$, $\operatorname{rk}_R(F)=n<\infty$, and $N\leq F$. Then $\exists m\leq n$ such that $N\cong R^m$.\\
Consequently, for $\mu_R(M)\defeq\min\{\abs{X}:\langle\!\langle X\rangle\!\rangle_R=M\}$, $\mu_R(M)\leq n$ implies $S\leq M\Longrightarrow \mu_R(S)\leq n$.
\begin{proof}
Identify $F\simeq R^n$. For $1\leq k\leq n$, set $I_k\defeq\{x_k:x=(x_1,\ldots,x_k,0,\ldots,0)\in N\}\lhd R$ and write $I_k=\langle r_k\rangle$..\\
If $r_k\neq0$, choose $u_k\in N\cap(R^k\times\{0\}^{n-k})$ with $(u_k)_k=r_k$. Let $K=\{k:r_k\neq0\}$.\\
If $\sum\limits_{k\in K}c_ku_k=0$ and $j=\max\{k:c_k\neq0\}$, then the $j$th coordinate gives $c_jr_j=0$, hence $c_j=0$, contradiction.\\
Thus $\{u_k:k\in K\}$ is $R$-independent. Set $\ell(0)=0$ and $\ell(x)=\max\operatorname{supp}(x)$ for $0\neq x\in N$.\\
If $j=\ell(x)>0$, then $x_j\in I_j=\langle r_j\rangle$, say $x_j=cr_j$, and $\ell(x-cu_j)<j$.\\
Induction on $\ell(x)$ gives $N=\bigoplus\limits_{k\in K}Ru_k\cong R^{\abs{K}}$, so $\operatorname{rk}_R(N)=\abs{K}\leq n$.\\
If $\mu_R(M)\leq n$, choose $\pi:R^n\twoheadrightarrow M$. For $S\leq M$, $P\defeq\pi^{-1}(S)\leq R^n$ has $P=\langle\!\langle v_1,\ldots,v_m\rangle\!\rangle_R$ with $m\leq n$.\\
Hence $S=\langle\!\langle \pi(v_1),\ldots,\pi(v_m)\rangle\!\rangle_R$ and $\mu_R(S)\leq m\leq n$.\\
\hlt{Structural Commentary}: PIDs let one choose pivot generators in a submodule of $R^n$.\\
\hlt{Key Insight}: Pivot ideals $I_k$ yield triangular free generators.
\end{proof}
\end{theorem}

\begin{corollary} \label{cor:finitely-generated-pid-modules-noetherian} \hlt{Finitely Generated PID Modules Are Noetherian}\\
If $\mu_R(M)<\infty$, then $M$ is Noetherian.
\begin{proof} Let $n=\mu_R(M)$. By Theorem \ref{thm:submodule-of-finite-free-pid-is-free}, $N\leq M\Longrightarrow \mu_R(N)\leq n<\infty$.\\
Thus every submodule of $M$ is f.g.; Theorem \ref{thm:noetherian-iff-submodules-finitely-generated} gives $M$ Noetherian.\\
\hlt{Structural Commentary}: Finite generation of every submodule prevents infinite strict growth.\\
\hlt{Key Insight}: Generators of the union appear at a finite stage.
\end{proof}
\end{corollary}

\begin{definition} \label{def:rank-finite-free-module-pid} \hlt{Rank of a Finite Free Module}\\
If $F$ is a finite free $R$-module and $F\cong R^n$, the integer $n$ is the \hlt{rank} of $F$, denoted $\operatorname{rk}(F)$.
\end{definition}

\begin{proposition} \label{prop:finite-free-rank-well-defined-domain} \hlt{Finite Free Rank Is Well-Defined over a Domain}\\
Let $R$ be an integral domain and $m,n\geq0$. Then $R^m\cong_R R^n\Longrightarrow m=n$.
\begin{proof}
Let $K=\operatorname{Frac}(R)$ and $\phi:R^m\xrightarrow{\sim}R^n$. Then $K\otimes_R\phi:K\otimes_R R^m\xrightarrow{\sim}K\otimes_R R^n$, i.e. $K^m\xrightarrow{\sim}K^n$.\\
Hence $m=\dim_KK^m=\dim_KK^n=n$.\\
\hlt{Structural Commentary}: Passing to the fraction field recovers ordinary dimension.\\
\hlt{Key Insight}: Rank over a domain is detected after denominators are allowed.
\end{proof}
\end{proposition}

\begin{proposition} \label{prop:spanning-set-size-rank-basis-pid} \hlt{A Full-Size Spanning Set Is a Basis}\\
Let $\operatorname{rk}_R(F)=n<\infty$. If $F=\langle\!\langle s_1,\ldots,s_n\rangle\!\rangle_R$, then $(s_1,\ldots,s_n)$ is an $R$-basis of $F$.
\begin{proof}
Choose a basis $(b_i)_{i=1}^n$ and define $\tau\in\operatorname{End}_R(F)$ by $\tau(b_i)=s_i$. Then $\tau$ is onto.\\
Since $F$ is free, $\exists\rho:F\to F$ with $\tau\rho=\operatorname{id}_F$; hence $F=\ker\tau\oplus\operatorname{im}\rho$ and $\operatorname{im}\rho\cong F$.\\
By Theorem \ref{thm:submodule-of-finite-free-pid-is-free}, $\ker\tau\cong R^r$ for some $r\leq n$. Thus $R^n\cong F\cong\ker\tau\oplus F\cong R^{r+n}$, so $r=0$ by Proposition
\ref{prop:finite-free-rank-well-defined-domain}. Hence $\ker\tau=0$, $\tau$ is an isomorphism, and $(s_i)=\tau(b_i)$ is a basis.\\
\hlt{Structural Commentary}: Surjective finite-free endomorphisms over a PID are injective.\\
\hlt{Key Insight}: Rank leaves no room for a split kernel.
\end{proof}
\end{proposition}

\begin{perspective} \label{per:smith-normal-form-and-module-presentations} \hlt{Smith Form as Diagonalized Relations}\\
For a presentation $R^n\xrightarrow{A}R^m\to M\to0$, Smith form replaces $A$ by a diagonal relation matrix, decoupling finite relations into cyclic constraints.
\end{perspective}

\begin{theorem} \label{thm:smith-normal-form-over-pid} \hlt{Smith Normal Form over a PID}\\
Let $A\in\operatorname{Mat}(m,n,R)$, viewed as $R^n\to R^m$.\\
$\exists$ invertible $P\in\operatorname{Mat}(m,m,R)$ and $Q\in\operatorname{Mat}(n,n,R)$ such that $PAQ=\operatorname{diag}(d_1,\ldots,d_r,0,\ldots,0)$,\\
where $d_i\in R\setminus\{0\}$, $d_1\mid\cdots\mid d_r$, and $r=\operatorname{rk}_R(\operatorname{im}A)$. The ideals $\langle d_i\rangle$ are unique.
\begin{proof}
\begin{enumerate}[label=(\roman*)]
\setlength{\itemsep}{0pt}
\item \hlt{Existence}: If $A=0$, done. Set $I(A)=\langle a_{ij}\rangle=\langle d\rangle$. Elementary operations move a Bezout combination equal to $d$ to position $(1,1)$. If a pivot $p\nmid a$, Bezout on $(p,a)$ replaces $p$ by $\gcd(p,a)$, so $\langle p\rangle$ strictly increases; ACC in a PID forces termination. Then $p$ divides every entry, so clear its row and column.\\
Induct on the complementary submatrix. Finally replace each bad adjacent pair $(a,b)$ by $(\gcd(a,b),\operatorname{lcm}(a,b))$ until $d_1\mid\cdots\mid d_r$.
\item \hlt{Uniqueness}: Let $\Delta_k(A)$ be the ideal of $k\times k$ minors and $\Delta_0(A)=R$. Matrix equivalence preserves each $\Delta_k$. In Smith form, $\Delta_k(A)=\langle d_1\cdots d_k\rangle$ for $1\leq k\leq r$, and $\Delta_k(A)=0$ for $k>r$. Since $R$ is a domain, $(\Delta_k(A):\Delta_{k-1}(A))=\langle d_k\rangle$, so $d_k$ is determined up to associates.
\end{enumerate}
\noindent
\hlt{Structural Commentary}: Smith form preserves rank plus divisibility data.\\
\hlt{Key Insight}: Determinantal ideals encode invariant factors.
\end{proof}
\end{theorem}

\begin{theorem} \label{thm:adapted-basis-submodule-free-pid} \hlt{Adapted Bases for Submodules of Free Modules}\\
Let $N\leq F$, $\operatorname{rk}_R(F)=n$, and $\operatorname{rk}_R(N)=k$.\\
Then $\exists(v_i)_{i=1}^n\in\operatorname{Bas}_R(F)$ and $r_1,\ldots,r_k\in R\setminus\{0\}$ such that $r_1\mid\cdots\mid r_k$ and $(r_iv_i)_{i=1}^k\in\operatorname{Bas}_R(N)$.
\begin{proof}
Choose bases of $N$ and $F$, and let $A$ represent $\iota:N\hookrightarrow F$.\\
Smith form gives $PAQ=\operatorname{diag}(r_1,\ldots,r_k,0,\ldots,0)$ with $r_i\neq0$ and $r_1\mid\cdots\mid r_k$, since $\iota$ is injective.\\
The row operations give a basis $(v_i)_{i=1}^n$ of $F$ and the column operations give a basis $(w_i)_{i=1}^k$ of $N$ with $\iota(w_i)=r_iv_i$. Hence $(r_iv_i)_{i=1}^k=(w_i)_{i=1}^k$ is a basis of $N$.\\
\hlt{Structural Commentary}: Inclusions diagonalize after choosing bases.\\
\hlt{Key Insight}: Diagonal entries are the required rescalings.
\end{proof}
\end{theorem}

\begin{corollary} \label{cor:smith-normal-form-presents-pid-modules} \hlt{Smith Form Presents Finitely Generated PID Modules}\\
Let $R^n\xrightarrow{A}R^m\to M\to0$ be a finite presentation.\\
If $A\sim\operatorname{diag}(d_1,\ldots,d_r,0,\ldots,0)$ in Smith form. $I=\{i:1\leq i\leq r,\ d_i\notin R^\times\}$, then $M\cong R^{m-r}\oplus\bigoplus_{i\in I}R/\langle d_i\rangle$.
\begin{proof}
Equivalence $A\sim PAQ$ changes bases in $R^m,R^n$, so $\operatorname{coker}A\cong\operatorname{coker}(PAQ)$.\\
For the diagonal map, $d_i\in R^\times$ gives $0$, $d_i\notin R^\times$ gives $R/\langle d_i\rangle$, and each zero row gives $R$.\\
\hlt{Structural Commentary}: Smith form turns presentations into direct sums.\\
\hlt{Key Insight}: Diagonal relations create cyclic torsion summands, and zero rows create free summands.
\end{proof}
\end{corollary}

\begin{method} \label{meth:smith-form-module-presentation} \hlt{Reading a Module from Smith Form}\\
For finite presentation $R^n\xrightarrow{A}R^m\to M\to0$ over PID, put $A\sim D=\operatorname{diag}(d_1,\ldots,d_r,0,\ldots,0)$, $d_1\mid\cdots\mid d_r$:
\begin{enumerate}[label=(\roman*)]
\setlength{\itemsep}{0pt}
\item $d_i\in R^\times\Longrightarrow$ no summand;
\item $d_i\notin R^\times\Longrightarrow R/\langle d_i\rangle$;
\item each zero target coordinate contributes one $R$;
\item factor nonunit $d_i$ into prime powers for elementary divisors.
\end{enumerate}
\hlt{Key Insight}: Smith form changes arbitrary relations into independent one-coordinate relations.
\end{method}

\begin{example} \label{ex:smith-normal-form-abelian-group-decomposition} \hlt{Smith Form for a Finite Abelian Group}\\
Let $M=\mathbb{Z}^2/\langle\!\langle(6,0),(0,10),(2,4)\rangle\!\rangle_{\mathbb{Z}}$, with relation matrix $A=\begin{bmatrix}6&0&2\\0&10&4\end{bmatrix}$. Then $\Delta_1(A)=\langle6,10,2,4\rangle=\langle2\rangle$ and $\Delta_2(A)=\langle60,24,-20\rangle=\langle4\rangle$. Hence $d_1\sim2$, $d_1d_2\sim4$, so $d_2\sim2$ and $M\cong\mathbb{Z}/2\mathbb{Z}\oplus\mathbb{Z}/2\mathbb{Z}$.\\
\hlt{Key Insight}: Determinantal ideals read off cyclic summands.
\end{example}

\begin{example} \label{ex:elementary-divisors-recombine-to-invariant-factors} \hlt{Recombining Elementary Divisors}\\
Let $G=\Z/12\Z\oplus\Z/18\Z$. Then $\operatorname{ElemDiv}(G)=\{2,2^2,3,3^2\}$. Recombining equal-height primary blocks gives $2\cdot3=6$ and $2^2\cdot3^2=36$; hence $G\cong\Z/6\Z\oplus\Z/36\Z$, with $6\mid36$.\\
\hlt{Structural Commentary}: Elementary divisors split by primes; invariant factors recombine them.\\
\hlt{Key Insight}: Prime-wise data equals one divisibility chain.
\end{example}
